\chapter{前向运动学 (Forward Kinematics)}

前向运动学描述机器人末端执行器（End-effector）位姿 $T \in SE(3)$ 与关节变量 $\theta \in \mathbb{R}^n$ 之间的函数关系：$T = f(\theta)$。

\section{指数积公式 (PoE Formula)}

指数积公式是 D-H 参数法的现代替代方案。其基本思想是：将机器人的运动看作是从一个“初始构型”开始，依次绕各个关节轴进行螺旋运动（Screw Motion）的过程。

\subsection{空间系指数积公式 (PoE in Space Frame)}



在空间系 PoE 中，所有的螺旋轴 $\mathcal{S}_i$ 都在固定基座坐标系 $\{s\}$ 中定义。

\begin{defbox}{零位构型 (Zero Configuration) $M$}
当所有关节变量 $\theta_1 = \theta_2 = \dots = \theta_n = 0$ 时，末端执行器坐标系 $\{b\}$ 相对于空间坐标系 $\{s\}$ 的位姿矩阵，记为 $M \in SE(3)$：
\[ M = T_{sb}(0) \]
\end{defbox}

\begin{formulabox}{空间系 PoE 公式}
对于一个 $n$ 轴开链机械臂，其前向运动学方程为：
\begin{equation}
T(\theta) = e^{[\mathcal{S}_1]\theta_1} e^{[\mathcal{S}_2]\theta_2} \cdots e^{[\mathcal{S}_n]\theta_n} M
\end{equation}
其中：
\begin{itemize}
    \item $\mathcal{S}_i = (\omega_i, v_i)$ 是当关节 $1, \dots, i-1$ 处于零位时，第 $i$ 个关节在空间系 $\{s\}$ 中的螺旋轴。
    \item $[\mathcal{S}_i] \in se(3)$ 是其矩阵表示。
\end{itemize}
\end{formulabox}

\textbf{推导逻辑：}
想象机器人从末端开始运动。首先，末端相对于基座的位置是 $M$。接着，绕轴 $n$ 运动 $\theta_n$ 产生变换 $e^{[\mathcal{S}_n]\theta_n}M$。以此类推，每一个后续的关节运动都相当于在其之前的运动链左侧施加一个新的变换。

\subsection{物体系指数积公式 (PoE in Body Frame)}

如果螺旋轴是在末端执行器当前的移动坐标系 $\{b\}$ 中定义的，则使用物体系 PoE。

\begin{formulabox}{物体系 PoE 公式}
\begin{equation}
T(\theta) = M e^{[\mathcal{B}_1]\theta_1} e^{[\mathcal{B}_2]\theta_2} \cdots e^{[\mathcal{B}_n]\theta_n}
\end{equation}
其中 $\mathcal{B}_i$ 是第 $i$ 个关节在物体系 $\{b\}$ 中的螺旋轴表达。
\end{formulabox}

\begin{mybox}{空间轴与物体轴的转换关系}
空间螺旋轴 $\mathcal{S}_i$ 与物体螺旋轴 $\mathcal{B}_i$ 之间可以通过零位构型 $M$ 的伴随变换相互转换：
\[ \mathcal{S}_i = [Ad_M] \mathcal{B}_i, \quad \mathcal{B}_i = [Ad_{M^{-1}}] \mathcal{S}_i \]
\end{mybox}

\section{螺旋轴 $\mathcal{S}_i$ 的确定方法}

在空间系中确定 $\mathcal{S}_i = (\omega_i, v_i)$ 的步骤如下：

\begin{enumerate}
    \item \textbf{转动关节 (Revolute Joint)}:
    \begin{itemize}
        \item $\omega_i \in \mathbb{R}^3$ 是关节轴 $i$ 的单位方向向量（遵循右手定则）。
        \item $v_i = -\omega_i \times q_i$，其中 $q_i$ 是关节轴 $i$ 上任意一点在 $\{s\}$ 中的坐标。
    \end{itemize}
    \item \textbf{移动关节 (Prismatic Joint)}:
    \begin{itemize}
        \item $\omega_i = 0$。
        \item $v_i \in \mathbb{R}^3$ 是关节平移方向的单位向量。
    \end{itemize}
\end{enumerate}



\section{PoE 与 D-H 参数法的对比}

\begin{table}[h]
\centering
\begin{tabular}{|l|p{6cm}|p{6cm}|}
\hline
\textbf{特性} & \textbf{指数积公式 (PoE)} & \textbf{D-H 参数法} \\ \hline
坐标系分配 & 仅需基座系 $\{s\}$ 和末端系 $\{b\}$ & 每个连杆必须按照特定规则建立坐标系 \\ \hline
参数含义 & 具有明确的几何意义（螺旋轴） & 较为抽象（偏置、扭角等） \\ \hline
奇异性 & 螺旋轴定义不具有局部奇异性 & 在平行轴等特殊排布下参数易突变 \\ \hline
计算结构 & 矩阵指数运算，数学形式统一 & 四个变换矩阵相乘 \\ \hline
\end{tabular}
\end{table}

\section{运动学总结与示例}

以 UR5 机器人为例，建立 PoE 模型的步骤：
1. **定义零位构型**：确定机器人所有关节为 0 状态下的 $M$ 矩阵。
2. **提取螺旋轴**：在 $M$ 状态下，观察每个关节轴线在 $\{s\}$ 系下的朝向 $\omega_i$ 和位置 $q_i$。
3. **代入 PoE 公式**：计算 $T(\theta)$。



\begin{mybox}{小贴士：下标规则}
在阅读本书公式时，注意 $e^{[\mathcal{S}]\theta}$ 始终作用于其右侧的元素。在空间系公式中，$M$ 在最右边，意味着变换是从“远端”向“基座”递归定义的；而在物体系公式中，$M$ 在最左边。
\end{mybox}

